\subsection{Scale}

All experiments are at $60$M--$300$M parameters. The trade-off
characterisation may shift at $\ge 1$B; published works at $7$B
\citep{zhou2024transfusion,gong2024diffullama} report different
direction signals. We treat the $60$M--$300$M characterisation as
"small-scale evidence" and explicitly do not extrapolate.

\subsection{Substrate}

All training is on GSM8K-train ($\sim 4$M tokens, math-reasoning
domain). The data-efficiency crossover is substrate-specific; we
make no claim that it generalises to broader corpora.

\subsection{Eval metric}

Downstream accuracy on GSM8K-dev is binomial-noise dominated at this
scale (\textless $6\%$ correct out of $50$ in most cells). We rely
on continuous NLL metrics throughout, but downstream-task accuracy
is the metric practitioners care about. Mode-switching at $n=8$
tightens the CI but does not eliminate the noise.

\subsection{Single hyperparameter trace}

The $\alpha$-schedule (linear $1.0 \to 0.5$), mask-ratio distribution
($\text{Uniform}(0.1, 0.9)$), and learning rate ($6\times 10^{-4}$
with cosine decay) were picked from BD3-LMs defaults and not swept.
A small $\alpha$-schedule sweep is the natural follow-up.

\subsection{Routing scope (Phase I narrow null)}

Phase~I (Section~\ref{sec:interleaved}) tested both scripted multi-round
AR$\leftrightarrow$diff interleaving and three non-learned heuristic
routers on F9 (FineWeb-only) and F10 across multiple training stages.
On F9 the diff-head loop pathology let fine-grained iteration look
like a win ($+2$\,pp accuracy, $-22$\,pp loop rate over the
single-switch baseline). On F10 at convergence the same iteration
ties pure AR ($4\,\%$ accuracy), and no heuristic crosses the
$+4$\,pp gate over the best fixed schedule. Per the Phase~I decision
rule we did not launch a learned router. The honest scope-limit is
specific: \emph{routing-during-AR-generation for accuracy lift} helps
when the diff head is broken, and disappears when proper training
(F10's Q/A mix) fixes the diff head directly. A learned router may
still earn its weight at scales where AR and diff heads diverge more
under joint training (Phase~I.2/I.3 future work).

This null sits alongside three positive Phase~J measurements
(Section~\ref{sec:phase-j}): the composite still delivers a
$5.43\times$ wall-clock speedup via parallel diff-fill, a FIM
capability AR structurally lacks, and a $-0.74$ NLL/token revision
improvement on gold continuation. The Phase~I null and the Phase~J
wins are independent claims about different axes; the null does not
undermine the wins.

\subsection{Interpretability controls}

The specialisation result (Section~\ref{sec:interp}) rests on a
best-cosine metric whose dynamic range is compressed in high dimension
(the chance floor is $0.124$, not $0$). We anchor the cross-mode
$0.169$ against a same-mode different-seed baseline ($0.329$, from a
matched-length $8{,}000$-step retrain), which is the correct null; a
shorter $2{,}000$-step pilot gave a more conservative $0.286$, and the
matched run widened the gap as expected. The null is measured at the
pre-head residual (\code{ln\_f})
only: the by-depth result (\textsection\ref{sec:p2b}) reports
cross-mode cosine at each of the ten blocks but not a per-block
same-mode null, so the depth-uniformity claim is anchored on the
\code{ln\_f} control rather than a null at every layer. The
causal-steering evidence (P.4) is \emph{suggestive, not
dispositive}: amplifying an AR-specific feature changes generation in
a content-conditional way, but we did not run the random-direction and
diff-feature steering controls that separate directed causation from
generic perturbation. We frame the steering as consistent-with, not
proof-of, the specialisation claim.

\subsection{Future work}

\begin{enumerate}
  \item Test the data-efficiency crossover on a non-math substrate
        (e.g., a $\sim 10$M-token slice of a code corpus).
  \item Add random-direction / diff-feature
        steering controls, to harden the specialisation result.
  \item Extend the compute-matched control to AR-only-at-$2\times$ as
        well, to test "does composite also beat compute-matched
        pure-AR on its own axis at small data?"
  \item Rerun Phase~I (interleaving + heuristic routers) on F10 at
        convergence (step $183$k, not yet available at submission).
        If the routing-vs-fixed picture shifts as the diff head
        specialises further, a REINFORCE-trained router becomes
        worth attempting; otherwise the negative reported here is the
        final word at $\sim$300\,M scale.
  \item Replicate at $1$B from scratch to test whether the crossover
        \emph{and} the routing-versus-fixed picture shift with scale.
\end{enumerate}
